% Part: lambda-calculus
% Chapter: lambda-definability
% Section: fixpoints

\documentclass[../../../include/open-logic-section]{subfiles}

\begin{document}

\olfileid[id]{lam}{ldf}{fp}

\olsection{Titik Tetap}

Andaikan kita ingin mendefinisikan fungsi faktorial secara rekursif sebagai
suatu suku $\fn{Fac}$ dengan sifat berikut:
\[
\fn{Fac} \ident \lambd[n][\fn{IsZero}\, n\, \num 1 (\fn{Mult}\, n (\fn{Fac}(\fn{Pred} \, n)))]
\]
Artinya, faktorial dari $n$ adalah $1$ jika $n = 0$, dan dalam kasus lainnya
merupakan $n$ dikalikan faktorial dari $n-1$. Tentu saja, kita tidak dapat
mendefinisikan suku $\fn{Fac}$ dengan cara ini karena $\fn{Fac}$ sendiri
muncul di ruas kanan. Definisi rekursif semacam itu, yang melibatkan acuan
diri, bukan bagian dari kalkulus lambda. Mendefinisikan suatu suku, misalnya,
dengan
\[
\fn{Mult} \ident \lambd[ab][a (\fn{Add}\, b) \num{0}]
\]
hanya melibatkan suku-suku yang telah didefinisikan sebelumnya di ruas kanan,
seperti $\fn{Add}$. Kita selalu dapat menyingkirkan $\fn{Add}$ dengan
menggantinya dengan suku pendefinisinya. Dengan demikian, suku $\fn{Mult}$
akan menjadi suku lambda murni; jika $\fn{Add}$ sendiri melibatkan suku-suku
terdefinisi (seperti halnya $\fn{Add}'$, misalnya), kita dapat meneruskan
proses ini dan akhirnya memperoleh suatu suku lambda murni.

Namun, hal ini tidak berlaku bagi definisi rekursif seperti definisi
$\fn{Fac}$ di atas. Jika kemunculan $\fn{Fac}$ di ruas kanan kita ganti dengan
definisi $\fn{Fac}$ itu sendiri, kita memperoleh:
\begin{align*}
  \fn{Fac} & \ident \lambd[n][\fn{IsZero}\, n\, \num{1}\,
    (\fn{Mult}\, n
  ((\lambd[n][\fn{IsZero} \, n \, \num 1\, (\fn{Mult}\, n\, (\fn{Fac}
    (\fn{Pred}\, n)))]) (\fn{Pred}\, n)))]
\end{align*}
dan $\fn{Fac}$ masih belum tersingkir dari ruas kanan. Jelas bahwa, jika kita
mengulangi proses ini, definisinya terus bertambah panjang dan proses tersebut
tidak pernah menghasilkan suku lambda murni. Jadi, cara ini tidak layak untuk
mendefinisikan faktorial (atau, secara lebih umum, fungsi rekursif).

Meskipun demikian, definisi rekursif tersebut memberi tahu kita sesuatu:
jika $f$ adalah suku yang merepresentasikan fungsi faktorial, maka suku
\[
\fn{Fac}' \ident \lambd[g][\lambd[n][\fn{IsZero} \, n \, \num 1\, (\fn{Mult} \, n\, (g (\fn{Pred} n)))]]
\]
yang diterapkan pada suku $f$, yakni $\fn{Fac}'\,f$, juga merepresentasikan
fungsi faktorial. Artinya, jika kita memandang $\fn{Fac}'$ sebagai fungsi yang
menerima suatu fungsi dan mengembalikan suatu fungsi, nilai
$\fn{Fac'}\, f$ tidak lain adalah~$f$, asalkan $f$ merupakan fungsi faktorial.
Suatu fungsi~$f$ dengan sifat $\fn{Fac}' \, f \equal[\beta] f$ disebut
\emph{titik tetap} dari $\fn{Fac}'$. Jadi, fungsi faktorial merupakan titik
tetap dari $\fn{Fac}'$.

Dalam kalkulus lambda terdapat suku-suku yang menghitung titik tetap dari suatu
suku yang diberikan, dan suku-suku tersebut kemudian dapat digunakan untuk
mengubah suku seperti $\fn{Fac}'$ menjadi definisi fungsi faktorial.

\begin{defn}
  \ollabel{defn:Turing-Y}

  \emph{Kombinator Y} adalah suku:
  \[
  Y \ident (\lambd[ux][x(uux)])(\lambd[ux][x(uux)]).
  \]
\end{defn}

\begin{thm}
  $Y$ mempunyai sifat bahwa $Yg \red g(Yg)$ untuk setiap suku $g$. Jadi,
  $Yg$ selalu merupakan titik tetap dari~$g$.
\end{thm}

\begin{proof}
  Mari singkat $(\lambd[ux][x(uux)])$ sebagai~$U$, sehingga $Y \ident UU$.
  Maka
  \begin{align*}
    Y g &\ident (\lambd[ux][x(uux)])U\,g \\
    &\red (\lambd[x][x(UUx)])g\\
    & \red g(UUg) \ident g(Yg).
  \end{align*}
  Karena $g(Yg)$ dan $Yg$ sama-sama tereduksi menjadi $g(Yg)$,
  $g(Yg) \equal[\beta] Yg$, sehingga $Yg$ merupakan titik tetap dari~$g$.
\end{proof}

Tentu saja, karena $Yg$ adalah suatu redex, reduksi dapat berlanjut tanpa
batas:
\begin{align*}
  Y g &\red g (Y g) \\
    &\red g (g (Y g)) \\
    &\red g(g (g (Y g)))\\
    &\ldots
\end{align*}
Jadi, kita dapat memandang $Yg$ sebagai $g$ yang diterapkan pada dirinya
sendiri tak berhingga kali. Jika kita menerapkan~$g$ satu kali lagi, maka,
boleh dikatakan, kita tidak melakukan hal tambahan apa pun; $g$ yang
diterapkan pada~$g$ yang diterapkan tak berhingga kali pada~$Yg$ tetaplah
$g$ yang diterapkan pada~$Yg$ tak berhingga kali.

Perhatikan bahwa barisan langkah reduksi-$\beta$ di atas yang dimulai
dengan~$Yg$ bersifat tak berhingga. Jadi, jika kita menerapkan $Yg$ pada suatu
suku, yakni meninjau $(Yg)N$, suku itu juga akan tereduksi menjadi tak
berhingga banyak suku yang berbeda, yaitu $(g(Yg))N$, $(g(g(Yg)))N$, \dots.
Meskipun demikian, mungkin saja suatu barisan langkah reduksi \emph{lain}
berakhir pada bentuk normal.

Sebagai contoh, tinjau fungsi faktorial. Definisikan $\fn{Fac}$ sebagai
$Y\, \fn{Fac}'$ (yakni, suatu titik tetap dari $\fn{Fac'}$). Maka:
\begin{align*}
  \fn{Fac}\,\num 3
  &\red Y\, \fn{Fac}' \, \num 3 \\
  &\red \fn{Fac}' (Y \,\fn{Fac}') \, \num 3 \\
  &\ident (\lambd[x][\lambd[n][\fn{IsZero} \, n\, \num 1\,
      (\fn{Mult}\, n\, (x (\fn{Pred}\, n)))]]) \, \fn{Fac} \, \num 3\\
  & \red \fn{IsZero} \, \num 3\, \num 1\,
      (\fn{Mult}\, \num 3\, (\fn{Fac} (\fn{Pred}\, \num 3))) \\
  &\red  \fn{Mult}\, \num 3 \, (\fn{Fac} \, \num 2).
  \intertext{Demikian pula,}
  \fn{Fac}\,\num 2
  &\red  \fn{Mult}\, \num 2 \, (\fn{Fac} \, \num 1) \\
  \fn{Fac}\,\num 1
  &\red  \fn{Mult}\, \num 1 \, (\fn{Fac} \, \num 0)
  \intertext{tetapi}
  \fn{Fac}\,\num 0
  &\red \fn{Fac}' (Y \,\fn{Fac}') \, \num 0 \\
  &\ident (\lambd[x][\lambd[n][\fn{IsZero} \, n\, \num 1\,
      (\fn{Mult}\, n\, (x (\fn{Pred}\, n)))]]) \, \fn{Fac} \, \num 0\\
  & \red \fn{IsZero} \, \num 0\, \num 1\,
      (\fn{Mult}\, \num 0\, (\fn{Fac} (\fn{Pred}\, \num 0))).\\
  &\red \num 1.
  \intertext{Jadi, secara keseluruhan,}
  \fn{Fac}\,\num 3
  &\red  \fn{Mult}\, \num 3 \,
  (\fn{Mult} \, \num 2\, (\fn{Mult} \, \num 1\, \num 1)).
\end{align*}

Apa yang berlaku bagi $\fn{Fac'}$ berlaku pula bagi setiap definisi rekursif.
Andaikan kita mempunyai persamaan rekursif
\begin{align*}
g\,x_1\dots x_n & \equal[\beta] N
\intertext{dengan $N$ mungkin memuat~$g$ dan $x_1$, \dots,~$x_n$. Maka selalu
terdapat suatu suku~$G \ident (Y \lambd[g][\lambd[x_1\dots x_n][N]])$
sedemikian sehingga}
G\,x_1 \dots x_n & \equal[\beta] \Subst{N}{G}{g}.
\intertext{Sebab, berdasarkan teorema titik tetap,}
G \ident (Y \lambd[g][\lambd[x_1\dots x_n][N]]) & \red \lambd[g][\lambd[x_1\dots x_n][N]](Y \lambd[g][\lambd[x_1\dots x_n][N]])\\
& \ident (\lambd[g][\lambd[x_1\dots x_n][N]])\,G
\intertext{dan akibatnya}
G\,x_1 \dots x_n
& \red (\lambd[g][\lambd[x_1\dots x_n][N]])\,G\,x_1\dots x_n\\
& \red (\lambd[x_1\dots x_n][\Subst{N}{G}{g}])\,x_1\dots x_n\\
  & \red \Subst{N}{G}{g}.
\end{align*}

Kombinator $Y$ pada \olref{defn:Turing-Y} dikemukakan oleh Alan Turing. Alonzo
Church telah mengusulkan versi lain yang akan kita sebut~$Y_C$:
\[
Y_C \ident \lambd[g][(\lambd[x][g(xx)])(\lambd[x][g(xx)])].
\]
Kombinator Church sedikit lebih lemah daripada kombinator Turing, dalam arti
bahwa $Y_Cg \equal[\beta] g(Y_Cg)$, tetapi reduksi satu-langkah
$Y_Cg \redone g(Y_Cg)$ tidak berlaku (meskipun keduanya tereduksi ke suku
yang sama). Misalkan $V$ adalah suku $\lambd[x][g(xx)]$, sehingga
$Y_C \ident \lambd[g][VV]$. Maka
\begin{align*}
  VV & \ident (\lambd[x][g(xx)])V \red g(VV)
  \text{ dan dengan demikian}\\
  Y_C g & \ident (\lambd[g][VV])g \red VV \red g(VV), \text{ tetapi juga}\\
  g(Y_C g) & \ident g((\lambd[g][VV])g) \red g(VV).
\end{align*}
Dengan kata lain, $Y_Cg$ dan $g(Y_Cg)$ tereduksi menjadi suku yang sama,
$g(VV)$; jadi $Y_Cg \equal[\beta] g(Y_Cg)$. Hal ini sering kali memadai untuk
penerapan.

\end{document}
